\section{Exact identification and the measured boundary}
\label{glob:section}\label{sec:global}

Proposition~\ref{mat:local} identifies the metric up to conformal
scale and identifies the normalized equations in harmonic frames
near a contact point.  We first show why, for a connection
Laplacian without an additional potential, this is already an
exact local isometry with a unitary connection identification.
We then return to the maximal match of Section~\ref{bdy:section}
and prove that it extends over the whole manifold and fixes every
part of the measured boundary.

\subsection{Removing the conformal factor}

\begin{lemma}[Exact local identification]\label{glob:conformal}
In the contact coordinates and unitary frames of
Proposition~\ref{bdy:contact}, suppose that on a connected
neighborhood \(\Omega\) of zero the metrics satisfy
\(g_2=e^{2\omega}g_1\), and that the equations in the harmonic
frames \(F_j\), after scalar normalization to the same principal
part, have equal first-order coefficients.  Then
\[
 g_2=g_1,\qquad J_{\mathrm{loc}}=F_1F_2^{-1}\in U(2),\qquad
 d_{A_1}(J_{\mathrm{loc}}u)=J_{\mathrm{loc}}d_{A_2}u
                         \quad\text{on }\Omega.
\]
Moreover this identification agrees with the old match on the
known side and matches all corresponding source evaluations
throughout \(\Omega\).
\end{lemma}

\begin{proof}
Let
\[
 C_j=F_j^{-1}(dF_j+A_jF_j).
\]
These are the connection matrices in the harmonic frames; they
need not be skew-Hermitian in the standard matrix inner product.
The negative of the actual connection Laplacian, expressed in
such a frame, is
\[
 \Delta_{g_j,C_j}
 =g_j^{il}\bigl((\partial_i+C_{j,i})(\partial_l+C_{j,l})
       -\Gamma^k_{il}(g_j)(\partial_k+C_{j,k})\bigr).
\]
Its zeroth-order term is zero, because the columns of \(F_j\)
are harmonic.  The common-principal-part hypothesis therefore
says
\begin{equation}\label{glob:equal-frame-operators}
             e^{2\omega}\Delta_{g_2,C_2}=\Delta_{g_1,C_1}.
\end{equation}
Indeed their second-order coefficients agree after the indicated
conformal normalization, their first-order coefficients agree
by assumption, and both zeroth-order coefficients vanish.

Put \(\phi=(n-2)\omega/2\).  The conformal change of the
Christoffel symbols gives
\[
 e^{2\omega}\Delta_{g_2,C_2}
 =\Delta_{g_1,C_2}
          +2\langle d\phi,\partial+C_2\rangle_{g_1}.
\]
Comparison of first-order terms with
\eqref{glob:equal-frame-operators} yields
\(C_2=C_1-d\phi\,I\).  Substitution, with matrix multiplication
in its displayed order, gives
\begin{align*}
 \Delta_{g_1,C_1-d\phi I}
       +2\langle d\phi,\partial+C_1-d\phi I\rangle_{g_1}
 &=\Delta_{g_1,C_1}
       -(\Delta_{g_1}\phi+|d\phi|_{g_1}^2)I.
\end{align*}
Here \(\Delta_{g_1}=\operatorname{div}_{g_1}\nabla_{g_1}\).
The zeroth-order terms in
\eqref{glob:equal-frame-operators} consequently imply
\[
 \Delta_{g_1}\phi+|d\phi|_{g_1}^2=0,
 \qquad \Delta_{g_1}(e^\phi)=0.
\]
On the known open side the metrics already agree, so
\(e^\phi=1\).  Unique continuation on \(\Omega\) gives
\(e^\phi=1\) everywhere.  Since \(n\ge3\) and \(\omega\) is
real, \(\omega=0\), and hence \(C_1=C_2\).

The last equality is equivalent to
\(d_{A_1}(F_1F_2^{-1}u)=F_1F_2^{-1}d_{A_2}u\).
Although the harmonic frames are not unitary,
\(J_{\mathrm{loc}}=F_1F_2^{-1}\) is.  To see this, parallel
transport the intertwining identity along any path starting in
the known side.  Both connection transports are unitary, and
\(J_{\mathrm{loc}}\) is the old unitary identification at the
initial point.  It remains unitary along the path; connectedness
of \(\Omega\) proves the assertion everywhere.

For an arbitrary source \(f\in C_c^\infty(U_0;\mathbb C^2)\),
both \(w_{1,f}\) and \(J_{\mathrm{loc}}w_{2,f}\), in the common
coordinate domain, are solutions of the first homogeneous
connection equation.  The charts avoid \(\overline{U_0}\).
Their difference vanishes on the known side, so unique
continuation makes it zero on \(\Omega\).  Thus all evaluation
identities hold, not only those for the finite family used in
the local argument.  On an overlap with the old match, the two
proposed first-side images of a point, together with their invertible fiber maps,
would impose an invertible linear relation between the two
first-side evaluation maps.  Joint surjectivity in
Lemma~\ref{bdy:evaluations} forces the image points to coincide.
Surjectivity at the second-side point then forces the fiber maps
to coincide as well.  Thus the local identification agrees with
the old match on every overlap.
\end{proof}

\subsection{The match fills both interiors}

The source-evaluation argument below follows the embedding architecture
of \cite[Section~6.1]{LLS20}.  The maximal matching, boundary completion,
and recovery on the whole measured patch adapt
\cite[Section~7]{ScalarCompanion2026}; we give the details for the
unitary bundle identification as well as the base map.

\begin{proposition}[Global interior identification]\label{glob:interior}
The maximal match of Lemma~\ref{bdy:matches} is defined on
\(N_2^\circ\), and its base map is an onto isometry
\(\Phi:N_2^\circ\to N_1^\circ\).
\end{proposition}

\begin{proof}
If its domain had an interior frontier,
Proposition~\ref{bdy:contact} and the dilation in
Section~\ref{bdy:dilation} would supply the local data used in
Sections~\ref{sph:section}--\ref{mat:section}.
Proposition~\ref{mat:local} gives conformally equal metrics and
equal normalized harmonic-frame equations on a neighborhood of
the contact point in the original charts.  Apply
Lemma~\ref{glob:conformal}.
Coordinate identity and \(J_{\mathrm{loc}}\) extend the old match
to a neighborhood containing the frontier point.  The concluding
assertion of Lemma~\ref{glob:conformal} ensures agreement on every
overlap.  This contradicts maximality.  The
domain, already nonempty and open, has no relative frontier in
the connected interior, and is therefore all of \(N_2^\circ\).

The image is open because \(\Phi\) is a local isometry.  To see
that it is relatively closed, suppose \(\Phi(p_k)\to q\in
N_1^\circ\).  Compactness gives a subsequence with \(p_k\to p\in
N_2\) and \(J(p_k)\to J_0\in U(2)\).  If \(p\in\partial N_2\),
continuity of the evaluations and \eqref{bdy:matching} would give
\(I_1(q)=J_0I_2(p)=0\), contrary to
Lemma~\ref{bdy:evaluations}.  Thus \(p\) is interior, and
continuity gives \(\Phi(p)=q\).  Connectedness of \(N_1^\circ\)
makes the image all of that interior.  Injectivity was proved in
Lemma~\ref{bdy:matches}; the inverse is smooth because the map
is a bijective local diffeomorphism.
\end{proof}

\subsection{Boundary extension and the original gauge}

\begin{proposition}[Completion and the measured patch]\label{glob:boundary}
The interior maps \((\Phi,J)\) extend smoothly to the extended
manifolds, with \(\Phi:N_2\to N_1\) a diffeomorphism and \(J\)
a unitary bundle map intertwining the connections.  They preserve
the original copies of \(M\) and satisfy
\[
                  \Phi|_\Gamma=\operatorname{Id},\qquad J|_\Gamma=I.
\]
\end{proposition}

\begin{proof}
An onto interior isometry preserves lengths of curves in both
directions and therefore preserves intrinsic distance.  The
metric completion of the interior of a compact smooth
Riemannian manifold with boundary is the manifold with its length
metric.  In particular, a curve touching the boundary can be
pushed slightly into an inward collar with arbitrarily small
change of length; uniform local metric comparisons identify
the resulting completion with the original compact topology.
Thus \(\Phi\) and its inverse extend to mutually inverse
homeomorphisms of \(N_2,N_1\), mapping boundary to boundary.

The extensions are smooth.  In a sufficiently thin uniform
collar, the distance \(s_j\) to \(\partial N_j\) is smooth and
its gradient generates inward unit normal flow.  The isometry
preserves distance to boundary, so
\(s_1\circ\Phi=s_2\), and in the interior it intertwines these
gradient fields.  On a fixed small level \(s_2=s_0>0\), the map
is already smooth.  For \(0\le s\le s_0\), express it using
that smooth level map and the smooth normal flows on the two
sides.  This gives a smooth extension to \(s=0\); the inverse
has the same construction.  Parallel transport for \(A_1,A_2\)
along these normal curves similarly extends \(J\) smoothly,
retaining unitarity and the connection intertwining identity.

The maps are identity on the open cap, hence on its closure by
continuity.  With closure taken in \(N_j\), the original interiors satisfy
\[
          M_j^\circ=N_j^\circ\setminus\overline E.
\]
Bijectivity and identity on \(\overline E\) preserve these
complements, and hence their closures, the original copies of
\(M\).  The restricted diffeomorphism fixes the attachment patch
\(P=\{b>0\}\), and \(J=I\) there.

We must recover the boundary identity on every part of
\(\Gamma\), which may have more than one component.  For
\(f\in C_c^\infty(\Gamma;\mathbb C^2)\), let \(u_j\) now denote
the harmonic solutions on the original manifolds with boundary
value \(f\).  On the second manifold define
\(\widetilde u_1=J^{-1}(u_1\circ\Phi)\).  The metric and
connection identities make \(\widetilde u_1\) second-harmonic.
On \(P\), it has the same boundary value as \(u_2\).  It also
has the same covariant normal derivative there.  Indeed equality
of the measured forms, tested against arbitrary functions
supported in \(P\), gives the equality of smooth densities
\[
 (\nabla^{A_1}_{\nu_1}u_1)\,dS_{g_1}
       =(\nabla^{A_2}_{\nu_2}u_2)\,dS_{g_2}
                                      \quad\text{on }P.
\]
The densities agree by Lemma~\ref{bdy:jets}, the isometry
transports the outward normal, and the bundle map equals \(I\)
on \(P\).  Thus \(\widetilde u_1-u_2\) has zero Cauchy data
there.  Boundary and interior unique continuation give
\(\widetilde u_1=u_2\) throughout \(M\).

Taking boundary traces yields
\begin{equation}\label{glob:boundary-tests}
             J(p)f(p)=f(\Phi(p))
 \quad\bigl(p\in\partial M,\ f\in C_c^\infty(\Gamma;\mathbb C^2)\bigr).
\end{equation}
If \(p\in\Gamma\) and \(\Phi(p)\ne p\), choose \(f\) nonzero
at \(p\), supported in a small neighborhood of \(p\) avoiding
\(\Phi(p)\).  Equation~\eqref{glob:boundary-tests} contradicts
invertibility of \(J(p)\).  Hence \(\Phi(p)=p\) on \(\Gamma\).
Arbitrary vector values of \(f(p)\) then give \(J(p)=I\) there.
\end{proof}

\begin{proof}[Proof of Theorem~\ref{thm:main}]
The boundary-identity gauges \(V_j\) of
Section~\ref{bdy:section} put the connections in normal gauge.
Lemma~\ref{bdy:jets} and Proposition~\ref{bdy:cap} produce common
exterior Green data.  Lemmas~\ref{bdy:evaluations} and
\ref{bdy:matches} define the maximal match, and
Proposition~\ref{bdy:contact} supplies the local contact data
wherever an interior frontier is present.  Proposition~\ref{mat:local}, followed by Lemma~\ref{glob:conformal}, removes
that frontier.  Propositions~\ref{glob:interior} and
\ref{glob:boundary} give a smooth global metric and connection
identification fixing the entire measured patch.

To distinguish the original connections in this final step, write
\(A_j^{\mathrm{orig}}\) for them, so the connections used in the
proof were
\[
 A_j=V_j^{-1}A_j^{\mathrm{orig}}V_j+V_j^{-1}dV_j.
\]
In the original trivializations the bundle map from second to
first fiber is
\[
                         U=(V_1\circ\Phi)J V_2^{-1}.
\]
It is smooth and unitary, and equals \(I\) on \(\Gamma\), since
both original gauges equal \(I\) on the entire boundary.
The intertwining identity is
\(d_{\Phi^*A_1^{\mathrm{orig}}}(Uu)=U d_{A_2^{\mathrm{orig}}}u\).
Expanding it gives
\[
 A_2^{\mathrm{orig}}
    =U^{-1}(\Phi^*A_1^{\mathrm{orig}})U+U^{-1}dU,
 \qquad g_2=\Phi^*g_1.
\]
These are the asserted identities.  The argument uses densities
and a single nonempty cap; it requires neither orientability nor
connectedness of the boundary or of \(\Gamma\).
\end{proof}
